\section*{Chapter \ref{ch:m1:pnl-and-positions} --- P\&L and the Accounting of a Position}

\begin{solution}{exo:m1:pnl-and-positions:1}
Gross $1.0 + 1.2 + 1.5 = \$3.7$ million, $1.85\times$ equity; net $2.2 - 1.5
= \$0.7$ million, $0.35\times$.
\end{solution}

\begin{solution}{exo:m1:pnl-and-positions:2}
After the first buy $\bar p = 20.00$; after the second $\bar p =
(500\times20.00 + 300\times20.40)/800 = 20.15$. The sale realises $600\times
0.35 = \$210$ and leaves 200 shares at $\bar p = 20.15$. Unrealised at 20.30:
$200 \times 0.15 = \$30$.
\end{solution}

\begin{solution}{exo:m1:pnl-and-positions:3}
Cash $-10\,000 - 6\,120 + 12\,300 = -\$3\,820$; $-3\,820 + 200\times20.30 =
\$240 = 210 + 30$.
\end{solution}

\begin{solution}{exo:m1:pnl-and-positions:4}
Short 1\,000 at 80.00; short 2\,000 at $\bar p = 80.50$. The purchase of
3\,000 closes the 2\,000, realising $2\,000\times(80.50-79.50) = \$2\,000$,
and opens a long of 1\,000 at 79.50. At 79.80 the unrealised is \$300 and the
total \$2\,300.
\end{solution}

\begin{solution}{exo:m1:pnl-and-positions:5}
Price: $50\,000\times(-10)\times0.80 = -\pounds400\,000$. Currency:
$50\,000\times200\times0.04 = +\pounds400\,000$. Cross:
$50\,000\times(-10)\times0.04 = -\pounds20\,000$. Total $-\pounds20\,000$:
the two large terms cancel and the cross term is the result.
\end{solution}

\begin{solution}{exo:m1:pnl-and-positions:6}
Carried: $-20\,000\times(29.70-30.00) = +\$6\,000$. New trades:
$8\,000\times0.10 + (-5\,000)\times(-0.20) = +\$1\,800$. Fees $-260$; borrow
$-120$; interest on proceeds $+65$. Total $+\$7\,485$. Closing \omterm{def:m1:pnl-and-positions:position}{position}
$-17\,000$.
\end{solution}

\begin{solution}{exo:m1:pnl-and-positions:7}
Keep a queue of signed lots; a fill of opposite sign consumes lots from the
front, realising $\pm n(p - p_{\text{lot}})$ with the sign of the lot, and
any remainder becomes a new lot. On exercise 4's fills: $1\,000\times0.50 +
1\,000\times1.50 = \$2\,000$, the same as \omterm{def:m1:pnl-and-positions:realised}{average cost}. The two conventions
agree whenever a \omterm{def:m1:pnl-and-positions:position}{position} is closed \emph{entirely}, since then every lot is
consumed; they differ only on partial closes.
\end{solution}

\begin{solution}{exo:m1:pnl-and-positions:8}
The desk closes its winners and keeps its losers open: each closed trade is
a gain, \omterm{def:m1:pnl-and-positions:realised}{realised P\&L} is positive every month, and the losses accumulate in
the unrealised line, which the report omits. Lot-level accounting with a free
choice of which lot to close makes this easiest, but any convention permits
it. \Cref{prop:m1:pnl-and-positions:identity} is immune: the \emph{total}
P\&L marked at an independent price, reported daily, shows the loss from the
first month. \omterm{def:m1:pnl-and-positions:realised}{Realised P\&L} and win rate are descriptions of when a trader
chose to close, not of how much was made.
\end{solution}

\begin{solution}{pb:m1:pnl-and-positions:1}
\textbf{1.} $12\,000 \to 5\,000 \to -4\,000 \to 2\,000 \to 12\,000$.
\textbf{2.} $592\,200 + 764\,100 - 505\,200 - 837\,000 = +\$14\,100$.
\textbf{3.} $32\,000 \times 0.0012 = \$38.40$.
\textbf{4.} Between the second and third fills, short 4\,000.
\textbf{5.} The \omterm{def:m1:pnl-and-positions:position}{position} is again 12\,000 shares, so the gross P\&L is
$14\,100 + 12\,000(m - 84.00)$. At 83.55: \$8\,700.
\textbf{6.} Mid 83.75: \$11\,100.
\textbf{7.} Auction 83.76: \$11\,220.
\textbf{8.} A long is closed at the bid, 83.70: \$10\,500.
\textbf{9.} The auction price: it is public, unique and cannot be influenced
by the trader's own last small print; the \$2\,520 between the highest and
lowest of these marks is otherwise hers to choose.
\textbf{10.} $12\,000\times(83.76-84.00) = -\$2\,880$.
\textbf{11.} $7\,000\times0.84 = 5\,880$; $9\,000\times1.14 = 10\,260$;
$6\,000\times(-0.44) = -2\,640$; $10\,000\times0.06 = 600$: $+\$14\,100$.
\textbf{12.} $-2\,880 + 14\,100 - 38.40 - 310 = \$10\,871.60$.
\textbf{13.} Gross $11\,220 = -2\,880 + 14\,100$: nothing is unexplained.
\textbf{14.} Yes, but look closer: nearly all the trading gain comes from
the two sales made above the close, that is, from having been \emph{less
long} while the price fell. Both lines are one decision --- the view that the
stock would fall --- and the second purchase gave part of it back.
\textbf{15.} Sell 7\,000: realised $7\,000\times2.10 = 14\,700$, long 5\,000 at
82.50. Sell 9\,000: closes 5\,000 ($+12\,000$, total 26\,700) and opens short
4\,000 at 84.90. Buy 6\,000: closes 4\,000 ($+2\,800$, total 29\,500) and opens
long 2\,000 at 84.20. Buy 10\,000: long 12\,000 at $\bar p = 83.783$.
\textbf{16.} $12\,000\times(83.76-83.783) = -\$280$.
\textbf{17.} $29\,500 - 280 = \$29\,220$ is the P\&L \emph{since the \omterm{def:m1:pnl-and-positions:position}{position}
was opened}. Subtract the \omterm{def:m1:pnl-and-positions:realised}{unrealised P\&L} carried in yesterday,
$12\,000\times(84.00-82.50) = \$18\,000$: \$11\,220.
\textbf{18.} A fill of 1\,000 shares missing from the desk's records (a
late or manually booked execution, or a fill for another desk's account
allocated here) --- check the broker's execution list against the fill log
by execution identifier. Then a corporate action or transfer processed by the
broker and not by the desk. Last, an error at the broker.
\textbf{19.} \textbf{\$10\,872}, net, at the official close: $+14\,100$ from
the day's trades, $-2\,880$ on the carried \omterm{def:m1:pnl-and-positions:position}{position}, $-348$ of fees and
financing.
\textbf{20.} One written marking policy --- official close for the daily
P\&L, net of all costs --- with every other figure labelled as an estimate.
\end{solution}

\begin{solution}{iq:m1:pnl-and-positions:1}
Realised is the gain on the part of the \omterm{def:m1:pnl-and-positions:position}{position} that has been closed,
measured against a cost convention; unrealised is the marked gain on what is
still open. The convention moves money between the two and between tax
years; it cannot change the total, which is cash plus \omterm{def:m1:pnl-and-positions:position}{position} times mark
and does not know in which order anything was bought.
\iqlookfor{the invariant stated as a formula.}
\end{solution}

\begin{solution}{iq:m1:pnl-and-positions:2}
Binary floating point cannot represent most decimal prices exactly; sums of
many fills accumulate errors that depend on the order of addition, so two
systems disagree by a cent and reconciliations break. Use integers in a
fixed minor unit (ticks, or ten-thousandths of the currency) for prices and
cash, a decimal type in reporting layers, and floating point only for derived
statistics such as an \omterm{def:m1:pnl-and-positions:realised}{average cost}.
\iqlookfor{order-dependence and reconciliation, not just ``rounding''.}
\end{solution}

\begin{solution}{iq:m1:pnl-and-positions:3}
At the mid, \$500\,000; to liquidate, the bid: \$499\,900. The last trade is
irrelevant: it is stale and away from the market. For two million shares no
screen price applies: the value is the mid less the expected cost of selling
that size, a square-root impact of some tens of basis points at least, and a
prudent mark carries that reserve.
\iqlookfor{ignoring the last print, and size-dependent valuation.}
\end{solution}

\begin{solution}{iq:m1:pnl-and-positions:4}
The broker's record is the legal one. Compare execution by execution, by
identifier, the broker's list with the strategy's fill log and the drop
copy: a fill missing on our side (lost acknowledgement, a partial fill after
a cancel request, a fill during a reconnect) is the usual cause. The system
should build its \omterm{def:m1:pnl-and-positions:position}{position} from the drop copy, reconcile it continuously
against its own belief, and on divergence stop quoting and alert: a strategy
trading on a wrong \omterm{def:m1:pnl-and-positions:position}{position} hedges the wrong way.
\iqlookfor{drop copy as the source of truth and ``stop trading'' as the
reflex.}
\end{solution}

\begin{solution}{iq:m1:pnl-and-positions:5}
A decomposition of the P\&L into causes --- carried \omterm{def:m1:pnl-and-positions:position}{positions} by risk
factor, new trades, \omterm{def:m1:pnl-and-positions:carry}{carry}, fees, currency --- that sums to the total.
Unexplained P\&L that is small and random is noise from approximations; a
persistent or growing one means the risk model does not describe the book (a
missing risk factor, wrong marks or stale parameters) or that trades are
missing or misbooked. Either way the desk does not know what it is exposed
to.
\iqlookfor{unexplained as a control signal, not an accounting nuisance.}
\end{solution}

\begin{solution}{iq:m1:pnl-and-positions:6}
Store the immutable event log: every fill, bust and correction as its own
event with the execution identifier, exchange timestamp, receive timestamp
and a reference to the event it amends. Derive everything else: quantity,
cash and fees are sums over the log, independent of order by
\cref{prop:m1:pnl-and-positions:identity}, so late and out-of-order fills
need no special case for the total. Order-dependent quantities (\omterm{def:m1:pnl-and-positions:realised}{average cost},
\omterm{def:m1:pnl-and-positions:realised}{realised P\&L}) are recomputed by replaying the log in exchange-time order
from the last checkpoint before the amended event. A bust is a negating
event, never a deletion, so that every published figure can be reproduced as
of any time.
\iqlookfor{event sourcing, and the observation that the exact total is
order-free.}
\end{solution}
